\documentclass[aspectratio=169,11pt]{beamer}

\usetheme{Madrid}
\usecolortheme{default}
\usefonttheme{professionalfonts}
\setbeamertemplate{navigation symbols}{}
\setbeamertemplate{footline}[frame number]

\usepackage{amsmath, amssymb, booktabs, graphicx}

\title{Labor Law, Enforcement, and State Capacity}
\subtitle{Labor Markets and Political/Cultural Institutions --- Week 2}
\author{Teaching deck draft}
\date{}

\begin{document}

\begin{frame}
  \titlepage
\end{frame}

\begin{frame}{Central question}
\begin{itemize}
    \item When do labor laws alter wages, employment, job quality, and worker protection?
    \item When do they remain rules on paper?
    \item Which margin moves: law, enforcement, knowledge, compliance, or incidence?
\end{itemize}
\end{frame}

\begin{frame}{Bridge from Week 1}
\begin{itemize}
    \item Week 1 defined institutions as labor-market mechanisms.
    \item Week 2 zooms in on formal law and the institutions that make it effective.
    \item The labor object is not doctrine by itself; it is wages, employment, formality, benefits, claims, and sorting.
    \item The formal--informal margin is part of the equilibrium object from the start.
\end{itemize}
\end{frame}

\begin{frame}{Law on the books vs. law in practice}
\[
L^{eff}_{jt}=L^{book}_{jt}\times E_{jt}\times K_{jt}
\]
\begin{itemize}
    \item $L^{book}_{jt}$: statutory or court-made rule.
    \item $E_{jt}$: enforcement, administrative reach, courts, sanctions, complaints.
    \item $K_{jt}$: worker and firm knowledge, salience, and practical ability to act.
    \item Teaching device: separate legal text from the rule workers and firms expect to bind.
\end{itemize}
\end{frame}

\begin{frame}{Six objects to keep separate}
\small
\begin{tabular}{p{0.29\linewidth} p{0.58\linewidth}}
\toprule
Object & Labor-economics interpretation \\
\midrule
Law on the books & formal statute, doctrine, mandate, or right \\
Effective law & expected constraint after implementation \\
Enforcement intensity & inspection, litigation, sanction, complaint credibility \\
Worker knowledge & awareness of rights, remedies, and retaliation risk \\
Realized compliance & actual wages, hours, benefits, procedures, records \\
Incidence & wages, employment, formality, benefits, sorting, welfare \\
\bottomrule
\end{tabular}
\end{frame}

\begin{frame}{Labor-law dimensions}
\small
\begin{tabular}{p{0.28\linewidth} p{0.34\linewidth} p{0.25\linewidth}}
\toprule
Dimension & Institutional object & Core labor margin \\
\midrule
Hiring and contract form & probation, written contracts, temporary work & access to formal jobs \\
Dismissal protection & firing restrictions, severance, wrongful discharge & separations, hiring, turnover \\
Wage and hours rules & minimum standards, overtime, payment rules & wages, hours, compliance \\
Mandated benefits & leave, pensions, payroll benefits & wage-benefit incidence \\
Anti-retaliation and disputes & complaints, remedies, adjudication & claims and bargaining \\
Inspection and sanctions & inspectors, penalties, repeat visits & compliance and sorting \\
\bottomrule
\end{tabular}
\end{frame}

\begin{frame}{Anchor papers: different treatments}
\small
\begin{tabular}{p{0.27\linewidth} p{0.28\linewidth} p{0.31\linewidth}}
\toprule
Anchor & Variation & Observed labor margin \\
\midrule
MacLeod & law and contract environment & credible employment promises \\
Autor--Donohue--Schwab & wrongful-discharge doctrines & employment and dismissal-sensitive margins \\
Almeida--Carneiro & inspection exposure & informality, wages, compliance \\
Bertrand--Crepon & information about labor law & beliefs and employment behavior \\
Amengual & inspector social linkages & enforcement pathways and compliance \\
Berliner et al. & state capacity and rights & broad labor-rights performance \\
\bottomrule
\end{tabular}
\end{frame}

\begin{frame}{Enforcement as institution}
\begin{itemize}
    \item Inspectors, courts, sanctions, complaints, and administrative records determine whether rules bind.
    \item Enforcement is costly and selective: agencies target visible firms, high-risk sectors, complaints, or politically salient workplaces.
    \item Selectivity changes incidence because covered and uncovered firms face different expected costs.
    \item Almeida--Carneiro: inspection intensity is a labor-market treatment, not just a background control.
\end{itemize}
\end{frame}

\begin{frame}{State capacity and heterogeneity}
\begin{itemize}
    \item Administrative reach: can the state find firms, workers, payrolls, hours, and benefits?
    \item Territorial heterogeneity: does the rule bind outside visible urban labor markets?
    \item Observability: are firms registered, records reliable, and violations measurable?
    \item Complementarity: unions, NGOs, and worker networks can help produce enforcement capacity.
\end{itemize}
\end{frame}

\begin{frame}{Worker knowledge and implementation}
\begin{itemize}
    \item Rights can fail if workers do not know the rule, remedies, or complaint channel.
    \item Rights can fail if workers fear retaliation or expect claims to be slow and costly.
    \item Firms can also misperceive legal constraints and adjust hiring or contract form.
    \item Bertrand--Crepon: information identifies knowledge and salience, not a new statute or stronger sanctions.
\end{itemize}
\end{frame}

\begin{frame}{Informality and equilibrium adjustment}
\small
\begin{tabular}{p{0.25\linewidth} p{0.30\linewidth} p{0.31\linewidth}}
\toprule
Margin & What enforcement changes & Interpretation problem \\
\midrule
Employment & hiring, separations, vacancies & direct effect may hide sorting \\
Wages & cash pay and wage offsets & benefits may be valued differently \\
Compliance & hours, benefits, records & observed compliance can coexist with evasion \\
Formality & registration and contract status & uncovered sector may absorb adjustment \\
Sorting & firm size, sector, contract type & treated composition may change \\
\bottomrule
\end{tabular}
\end{frame}

\begin{frame}{Empirical designs and what they identify}
\small
\begin{tabular}{p{0.32\linewidth} p{0.28\linewidth} p{0.27\linewidth}}
\toprule
Design & Identifies best & Typical outcome \\
\midrule
Legal adoption or repeal & law on the books & employment, wages, turnover \\
Court-made variation & contract doctrine & dismissal-sensitive margins \\
Inspector access/allocation & enforcement intensity & compliance, formality, wages \\
Worker-information RCTs & knowledge and salience & beliefs, claims, employment \\
Institutional indices & broad legal environment & comparative rights and labor outcomes \\
Employer--employee admin data & incidence and sorting & wages, benefits, flows \\
\bottomrule
\end{tabular}
\end{frame}

\begin{frame}{What not to do}
\begin{itemize}
    \item Do not call a statute the treatment if the design varies enforcement.
    \item Do not call enforcement an outcome if it is the identifying variation.
    \item Do not interpret compliance without naming the covered and uncovered margins.
    \item Do not read employment effects without asking about wages, benefits, formality, and sorting.
\end{itemize}
\end{frame}

\begin{frame}{Week 2 lab}
\begin{itemize}
    \item Reproduce: synthetic Almeida--Carneiro style enforcement and informality panel.
    \item Diagnose: compare enforcement variation with Autor--Donohue--Schwab legal variation.
    \item Transfer: classify whether a design identifies law, enforcement, knowledge, or effective compliance.
    \item The smoke path is local and synthetic; the purpose is design diagnosis.
\end{itemize}
\end{frame}

\begin{frame}{Bridge to Week 3}
\begin{itemize}
    \item Week 2: when does formal labor law become effective?
    \item Week 3: how are labor markets organized when formal contracts and protections cover only part of work?
    \item The connecting object is the formal--informal margin.
    \item Keep separating legal protection, enforcement, knowledge, compliance, and incidence.
\end{itemize}
\end{frame}

\begin{frame}{Takeaways}
\begin{itemize}
    \item Labor law is a labor-market institution because it changes contracts, compliance, bargaining, and adjustment.
    \item Effective law depends on legal text, enforcement, and knowledge.
    \item State capacity and worker organization shape where law binds.
    \item Strong empirical work names the variation, observed margin, and object identified.
\end{itemize}
\end{frame}

\end{document}
